% SPDX-FileCopyrightText: Copyright (c) 2016-2026 Objectionary.com
% SPDX-License-Identifier: MIT

\begin{proof}[Proof of \cref{thm:morphing-confluent}]
Unlike normalization, whose rules may rewrite a redex at any position, morphing
  acts only at the root of its argument \(n\): every rule of \cref{fig:morphing}
  matches the whole of \(n\) against \(e\) and recurses through its premises on
  determined subterms.
It therefore suffices to show the rules are deterministic---at most one applies to
  a pair \((n, e)\), and the one that applies fixes the result---for then the
  relation they define is single-valued, which is the claim.
Determinism neither needs nor implies termination, so this holds even though
  morphing is partial (\cref{sec:morphing}).

\emph{At most one rule matches.}
By the grammar of \cref{fig:ebnf}, \(n\) is a locator ($Q$ or \(\phiTerminal{\xi}\)),
  the terminator $T$, a formation, a dispatch, or an application, and its outermost
  constructor already separates most rules.
For $T$ only \nameref{r:dead} matches and for \(\phiTerminal{\xi}\) only
  \nameref{r:xi}.
For $Q$ the universe decides: \nameref{r:universe} requires \(e \neq Q\) and
  \nameref{r:mg} requires \(e = Q\), so exactly one fits a given \(e\).
A formation \([[ B ]]\) with no head operation matches only \nameref{r:mf}.
A dispatch \(m.\tau\) is shared by \nameref{r:ml}, \nameref{r:mphi}, and
  \nameref{r:md}, which their conditions render disjoint: \nameref{r:md} demands a
  non-formation subject (\(\phinoNotFormation{m}\)), whereas \nameref{r:ml} and
  \nameref{r:mphi} take a formation one, and among those \nameref{r:ml} requires an
  $L$-asset while \nameref{r:mphi} forbids one (\(L \notin B\)).
The formation-headed dispatches left over---a subject with neither an $L$-asset nor
  an $@$-attribute, or one already binding \(\tau\)---match no rule, but they are
  never normal forms: normalization reduces them by \nameref{r:stop} or
  \nameref{r:dot} (\cref{sec:normalization}), and morphing is only ever applied to a
  normal form, so it never meets them.
An application splits four ways on its argument, an exhaustive and disjoint
  partition: named or positional (\(\tau -> a\) versus
  \(\phiTerminal{\alpha_i} -> a\)) crossed with absolute or not (\(a \in \mathcal{K}\)
  for \nameref{r:ma} and \nameref{r:maa}, \(a \notin \mathcal{K}\) for
  \nameref{r:mad} and \nameref{r:maad}), and every argument falls in one class.

\emph{The matching rule fixes the result.}
Its meta-variables are pinned by the term: a formation holds at most one $L$-asset,
  so \nameref{r:ml} splits \([[ B_1, L> F, B_2 ]]\) uniquely, and the subject and
  attribute of a dispatch, and the subject and argument of an application, are read
  off the syntax.
Each premise is then a single-valued judgment.
A normalization premise yields the unique normal form guaranteed by
  \cref{thm:normalization-total,thm:normalization-confluent}.
The evaluation premise of \nameref{r:ml} calls the function attached to the
  $L$-asset, which by \cref{sec:evaluation} maps a formation, the universe, and a
  state deterministically to an expression and a state, and then normalizes that
  expression to a unique normal form.
Every morphing premise is a shorter derivation, single-valued by the induction
  hypothesis.
So the rule maps its configuration, state included, to one result.

By induction on the derivation, then, \(\phinoMorph{n}{e}{s_1}{n_2}{s_2}\)
  determines \(n_2\) and \(s_2\) from \(n\), \(e\), and \(s_1\): a normal form morphs
  to at most one formation or $T$.
\end{proof}
